\documentclass[12font]{article}
\title{\textbf{ Note -- a Special Case for $N^A=3$ }}
\author{Chen-Chih Wang  }
\date{\textit{Department of physics, National Tsing Hua University, Hsinchu 30013, Taiwan} \\ (Dated: \today)}
\usepackage[margin=2cm]{geometry}
%\usepackage[fleqn]{amsmath}
%\usepackage{CJK,amsmath}
\usepackage{amsfonts,amssymb}
\usepackage{fancyhdr}
%\pagestyle{fancy}
\usepackage{textcomp}
\usepackage{graphicx}
\usepackage{float}
\usepackage{color}
\usepackage{tikz}
\usepackage{multicol}
\usepackage{threeparttable}
\usepackage{amsthm}
\usepackage{mathtools}
\usepackage{hyperref}
\urlstyle{same}
\usepackage{caption}
\usepackage{subcaption}
\usepackage{dsfont}
\usepackage{comment}
\newtheorem{theorem}{Theorem}[section]
\theoremstyle{definition}
\newtheorem{definition}{Definition}[section]
\newtheorem{lemma}[theorem]{Lemma}
\newtheorem{corollary}{Corollary}[theorem]
\newtheorem{proposition}[theorem]{Proposition}

\usepackage{cancel}

%\numberwithin{equation}{section}

\newcommand{\vac}{\mrm{vac}}
\newcommand{\Bra}[1]{\left\langle #1 \right|}
\newcommand{\Ket}[1]{\left| #1 \right\rangle}
\newcommand{\bra}[1]{\langle #1 |}
\newcommand{\ket}[1]{| #1 \rangle}
\newcommand{\anglelr}[1]{\left\langle #1 \right\rangle}
\newcommand{\sanglelr}[1]{\langle #1 \rangle}
\newcommand{\braket}[2]{\langle #1 | #2 \rangle}
\newcommand{\Bbraket}[2]{\left\langle #1 \right|\left. #2 \right\rangle}
\newcommand{\dif}[1]{\frac{\mathrm{d}}{\mathrm{d}#1}}
\newcommand{\diff}[2]{\frac{\mathrm{d}#1}{\mathrm{d}#2}}
\newcommand{\gd}{\boldsymbol{\nabla}}
\newcommand{\dive}{\boldsymbol{\nabla}\cdot}
\newcommand{\curl}{\boldsymbol{\nabla}\times}
\newcommand{\lr}[1]{\left(  #1  \right)}
\newcommand{\lrm}[1]{\left[  #1  \right]}
\newcommand{\lrg}[1]{\left\{  #1  \right\}}
\newcommand{\abs}[1]{\left|  #1  \right|}
\newcommand{\de}{\partial}
\newcommand{\pdif}[1]{\frac{\partial}{\partial#1}}
\newcommand{\pdiff}[2]{\frac{\partial#1}{\partial#2}}
\newcommand{\mrm}[1]{\mathrm{#1}}
\newcommand{\bo}[1]{\mathbf{#1}}
\newcommand{\bog}[1]{\boldsymbol{#1}}
\newcommand{\vint}[3]{\int_{#1}#2\cdot\mathrm{d}\mathbf{#3}}
\newcommand{\voint}[3]{\oint_{#1}#2\cdot\mathrm{d}\mathbf{#3}}
\newcommand{\Vint}{\int\mrm{d}^3\bo{r}'}
\newcommand{\rr}{\mathbf{r} - \mathbf{r}'}
\newcommand{\arr}{\left|\mathbf{r} - \mathbf{r}'\right|}
\newcommand{\nx}{\hat{\bo{n}}_1}
\newcommand{\ny}{\hat{\bo{n}}_2}
\newcommand{\alert}[1]{{\color{red}#1}}
\newcommand{\alertV}[1]{{\color{violet}#1}}
\newcommand{\Tr}{\mrm{Tr}}
\newcommand{\tr}{\mrm{tr}}
\newcommand{\norm}[1]{\left\lVert#1\right\rVert}

\newcommand{\Even}{\bo{Even}}
\newcommand{\Odd}{\bo{Odd}}
\newcommand{\TReven}{\bo{TReven}}
\newcommand{\TRodd}{\bo{TRodd}}
\newcommand{\Good}{\bo{G}}

\newcommand{\wordfig}[2]{\begin{minipage}[h]{#1 cm}
	\vspace{0pt}
	\includegraphics[width=\textwidth]{#2}
   \end{minipage}}

\begin{document}
\renewcommand\qedsymbol{$\blacksquare$}
\maketitle
%\fontsize{12pt}{18pt}\selectfont
According to the Yao-Qi method, the entanglement entropy of a subregion $A$ of the Kitaev model can be expressed as
\begin{equation}\label{e.Yao-Qi}
  S_A = -\tr\lr{\frac{\Gamma_A}{2}\ln\frac{\Gamma_A}{2}} + \lr{\frac{1}{2}\abs{\de A} - 1}\ln 2
\end{equation}
where $\Gamma_A$ is a $N^A\times N^A$ matrix defined by
\[
[\Gamma^A]_{ij} = \anglelr{c_ic_j}\quad;\quad i,j\in A
\]
If we take $N^A=3$, and define $\anglelr{c_1c_2}=ix$, $\anglelr{c_3c_2}=iy$, and $\eta=\sqrt{x^2+y^2}$, then the matrix $\Gamma_A$ reads
\[
\Gamma_A = \begin{pmatrix}
             1 & ix & 0 \\
             -ix & 1 & -iy \\
             0 & iy & 1 
           \end{pmatrix}
\]
The three eigenvalues of such a matrix are $\lrg{1,1\pm\eta}$. We also know that the length of the circumference of $A$ is $\abs{\de A}=5$. Thus,  equation (\ref{e.Yao-Qi}) gives the entanglement entropy as following
\begin{equation}\label{e.result-Yao-Qi}
S_A = 3\ln2 - \frac{1}{2}\lrm{\lr{1+\eta}\ln\lr{1+\eta} + \lr{1-\eta}\ln\lr{1-\eta}}
\end{equation}

On the other hand, thanks to the small size of $A$ we chose, the entanglement entropy can be calculated via exactly diagonalizing the reduced density matrix. The Majorana fermion representation of the reduced density matrix is
\[
\rho^F_A = \frac{1}{4}\lr{1 + \anglelr{ic_1c_2}ic_1c_2 + \anglelr{ic_3c_2}ic_3c_2} = \frac{1}{4}\lr{1 - xic_1c_2 - yic_3c_2}
\]
The factor $\frac{1}{4}$ corresponds to $2^{-\lceil N^A/2\rceil}$. Our next goal is to rearrange the formula into the Gibbs form $\mathcal{K}e^{-H_\mrm{ent}}$ to get the entanglement Hamiltonian $H_{\mrm{ent}}$. In this way, we can obtain the full information of the entanglement spectrum, and naturally get the value of the entanglement entropy. To do this, we define $\gamma_1=\frac{x}{\eta}c_1 + \frac{y}{\eta}c_3$ and $\gamma_2=c_2$. After that, the Majorana reduced density matrix becomes
\[
\rho^F_A = \frac{1}{4}\lrm{1 + \eta \lr{-i\gamma_1\gamma_2}} = \frac{1}{4\cosh\epsilon}\lrm{\cosh\epsilon+\sinh\epsilon\lr{-i\gamma_1\gamma_2}} = \frac{1}{4\cosh\epsilon}\exp\lr{-i\epsilon\gamma_1\gamma_2}
\]
where we have defined $\epsilon$ in the way such that $\tanh\epsilon=\eta$. The corresponding entanglement entropy is then
\begin{equation}\label{e.ED-first}
S^F_A = -\Tr_F\lr{\rho^F_A\ln\rho^F_A} = \epsilon\anglelr{i\gamma_1\gamma_2} + \ln 4 + \ln \cosh\epsilon
\end{equation}
Represent $\gamma_1$ and $\gamma_2$ in terms of $c$-Majorana fermions, we get
\[
\anglelr{i\gamma_1\gamma_2} = \frac{x\anglelr{ic_1c_2} + y\anglelr{ic_3c_2}}{\eta} = -\frac{x^2+y^2}{\sqrt{x^2+y^2}} = -\eta
\]
On the other hand, the term $\ln\cosh\epsilon$ can also be simplified,
\[
\ln\cosh\epsilon = -\frac{1}{2}\ln\lr{1-\tanh^2\epsilon} = -\frac{1}{2}\lrm{\ln\lr{1+\eta} + \ln\lr{1-\eta}}
\]
At the same time, $\epsilon$ can also be represented in terms of $\eta$ in a simple way,
\[
\epsilon = \tanh^{-1}\eta = \frac{1}{2}\lrm{\ln\lr{1+\eta} - \ln\lr{1-\eta}}
\]
Put all of these expressions back to (\ref{e.ED-first}), we obtain
\begin{equation}\label{e.ED-Majo}
S^F_A = 2\ln2 - \frac{1}{2}\lrm{\lr{1+\eta}\ln\lr{1+\eta} + \lr{1-\eta}\ln\lr{1-\eta}}
\end{equation}
What we have obtained is merely the fermion part of the entanglement entropy. To obtain the correct value of the entanglement entropy, we have to add the corresponding gauge contribution back, which is
\[
S^G_A = \lr{N^A-\left\lceil \frac{N^A}{2} \right\rceil - N^A_p}\ln 2 = \ln2
\]
So, eventually, we get
\begin{equation}\label{e.ED-result}
  S_A = S^F_A + S^G_A = 3\ln2 - \frac{1}{2}\lrm{\lr{1+\eta}\ln\lr{1+\eta} + \lr{1-\eta}\ln\lr{1-\eta}}
\end{equation}
which reproduces the value obtained by the Yao-Qi formula (\ref{e.result-Yao-Qi}). 

The gauge contribution $S^G_A$ arises from the degeneracy of the entanglement spectrum
\[
\rho_A = \frac{1}{2^{N^A}}\lr{\mathds{1}+\anglelr{\sigma^x_1\sigma^x_2}\sigma^x_1\sigma^x_2 + \anglelr{\sigma^y_3\sigma^y_2}\sigma^y_3\sigma^y_2} \cong \frac{1}{2^{N^A-\lceil N^A/2\rceil}}\begin{pmatrix}
                                      \rho^F_A & \; \\
                                      \; & \rho^F_A 
                                    \end{pmatrix}
\]
We can see the factor $\frac{1}{2^{N^A-\lceil N^A/2\rceil}}$ counts the degeneracy, and which contributes to $S^G_A$ we have obtained. 

To confirm this, we can again diagonalize the density matrix $\rho_A$. Such a density matrix can be written as
\[
\rho_A = \frac{1}{8}\lrm{\mathds{1} - \eta\lr{\frac{x}{\eta}\sigma^x_1\sigma^x_2 + \frac{y}{\eta}\sigma^y_3\sigma^y_2}} = \frac{1}{8\cosh\epsilon}\lrm{\cosh\epsilon\mathds{1} - \sinh\epsilon\lr{\frac{x}{\eta}\sigma^x_1\sigma^x_2 + \frac{y}{\eta}\sigma^y_3\sigma^y_2}}
\]
and one can observe that the operator $\lr{\frac{x}{\eta}\sigma^x_1\sigma^x_2 + \frac{y}{\eta}\sigma^y_3\sigma^y_2}$ is an $\mathbb{Z}_2$ operator. 
\[
\lr{\frac{x}{\eta}\sigma^x_1\sigma^x_2 + \frac{y}{\eta}\sigma^y_3\sigma^y_2}^2 = \frac{1}{\eta^2}\lr{x^2\mathds{1} + y^2\mathds{1}} = \mathds{1}
\]
So the density matrix can be written as the Gibbs form
\[
    \rho_A = \frac{1}{8\cosh\epsilon}\exp\lrm{-\epsilon\lr{\frac{x}{\eta}\sigma^x_1\sigma^x_2 + \frac{y}{\eta}\sigma^y_3\sigma^y_2}}
\]


One should notice that there's no $N_p^A$ in the formula because $A$ is too small so that there's no plaquette exclusively inside $A$, thus $N^A_p=0$ in the case of $N^A=3$. If we consider larger subregions, we have to include the projector $\mathcal{P}_A=\prod_{p\in A}\lr{\frac{\mathds{1}+W_p}{2}}$ inside the reduced density matrix $\rho_A$, and the factor $N^A_p$ will appear naturally.


\end{document}
To compare the spin representation with the Majorana representation, we take the transformation \[\sigma^{\alpha_{ij}}_i\sigma^{\alpha_{ij}}_j = \mathcal{P}_G\hat{u}_{ij}\lr{-ic_ic_j}\mathcal{P}_G\] 
to represent the spin density matrix, then we have
\[
\rho_A = \frac{1}{2}\mathcal{P}_G\lr{\mathds{1}+\hat{u}_{12}} \lr{\mathds{1}+\hat{u}_{32}}\rho^F_A\mathcal{P}_G
\] 